\documentclass[14pt,a4paper,oneside]{extbook}
\usepackage[utf8]{vietnam}
\usepackage{amsmath,amssymb,mathrsfs,fancyhdr,enumerate,multirow,makecell}
\usepackage{tkz-euclide,tikz-3dplot,tikz,tkz-tab}
\usetikzlibrary{arrows.meta,calc,intersections,angles,quotes,backgrounds,shapes.geometric,patterns,through,shadows,positioning}
\usepackage{yhmath}
\usepackage[top=0.7cm, bottom=0.7cm, left=1.0cm, right=1.0cm, headheight=16pt, headsep=3mm, footskip=5mm]{geometry}
\usepackage[hidelinks,unicode]{hyperref}
\usepackage{qrcode}
\usepackage[many]{tcolorbox}
\usepackage{fontawesome}

\pagestyle{fancy}
\fancyhf{}
\fancyhead[L]{\footnotesize\bfseries\color{teal!85!black}\faBook~SỔ TAY KIẾN THỨC TOÁN 4}
\fancyhead[R]{\footnotesize\bfseries\color{blue!85!black}Lớp Toán Thầy Lĩnh -- loptoanthaylinh.vn}
\fancyfoot[C]{\footnotesize\bfseries\color{gray!80!black}-- Sổ tay kiến thức: Bài 15. Làm quen với dãy số tự nhiên --}
\renewcommand{\headrulewidth}{0.4pt}

\begin{document}

% --- BANNER TIÊU ĐỀ BÀI HỌC ---
\noindent
\begin{tcolorbox}[
  enhanced,
  colback=teal!8!white,
  colframe=teal!80!black,
  arc=3mm, boxrule=1pt,
  left=3mm, right=3mm, top=2mm, bottom=2mm,
  halign=center
]
  {\Large\bfseries\color{teal!90!black}\faBookmark\ BÀI 15: LÀM QUEN VỚI DÃY SỐ TỰ NHIÊN}\\[1.5pt]
  {\footnotesize\itshape\color{black!75} Đặc điểm dãy số tự nhiên, tia số, số chẵn, số lẻ và hai số tự nhiên liên tiếp trên 1 trang A4}
\end{tcolorbox}

\vspace{2mm}

% --- KHỐI 1 & 2: LÝ THUYẾT & KĨ NĂNG ---
\noindent
\begin{minipage}[t]{0.49\textwidth}
  \begin{tcolorbox}[
    enhanced,
    title={\footnotesize\bfseries\color{white}\faListUl\ 1. Dãy số tự nhiên \& Đặc điểm},
    colback=blue!2!white,
    colframe=blue!75!black,
    arc=2.5mm, boxrule=0.8pt, top=1.8mm, bottom=1.8mm, left=2mm, right=2mm
  ]
    \footnotesize
    $\bullet$ Các số $0, 1, 2, 3, 4, 5, ...$ gọi là các \textbf{số tự nhiên}. Xếp theo thứ tự từ bé đến lớn tạo thành \textbf{dãy số tự nhiên}:
    \begin{center}
      $0,\, 1,\, 2,\, 3,\, 4,\, 5,\, 6,\, 7,\, 8,\, 9,\, 10,\, ...$
    \end{center}
    $\bullet$ \textbf{Số bé nhất} của dãy số tự nhiên là số \textbf{0}.\\
    $\bullet$ \textbf{Không có số tự nhiên lớn nhất}, dãy số tự nhiên kéo dài vô tận.\\
    $\bullet$ Thêm 1 vào bất kì số nào ta được số tự nhiên liền sau: $n + 1$. Bớt 1 ở số tự nhiên khác 0 ta được số liền trước: $n - 1$.
  \end{tcolorbox}
\end{minipage}%
\hfill
\begin{minipage}[t]{0.49\textwidth}
  \begin{tcolorbox}[
    enhanced,
    title={\footnotesize\bfseries\color{white}\faArrowsH\ 2. Tia số, Số chẵn \& Số lẻ},
    colback=teal!2!white,
    colframe=teal!75!black,
    arc=2.5mm, boxrule=0.8pt, top=1.8mm, bottom=1.8mm, left=2mm, right=2mm
  ]
    \footnotesize
    $\bullet$ \textbf{Biểu diễn trên tia số:}\\
    Gốc tia số là điểm 0. Chiều mũi tên từ trái sang phải biểu thị số lớn dần.\\
    $\bullet$ \textbf{Hai số tự nhiên liên tiếp}: Hơn kém nhau \textbf{1 đơn vị}. Trong hai số tự nhiên liên tiếp luôn có một số chẵn và một số lẻ.\\
    $\bullet$ \textbf{Số chẵn}: Có chữ số tận cùng là $0, 2, 4, 6, 8$. Hai số chẵn liên tiếp hơn kém nhau \textbf{2 đơn vị}.\\
    $\bullet$ \textbf{Số lẻ}: Có chữ số tận cùng là $1, 3, 5, 7, 9$. Hai số lẻ liên tiếp hơn kém nhau \textbf{2 đơn vị}.
  \end{tcolorbox}
\end{minipage}

\vspace{2mm}

% --- KHỐI 3: SƠ ĐỒ TƯ DUY TIKZ ---
\noindent
\begin{tcolorbox}[
  enhanced,
  title={\footnotesize\bfseries\color{white}\faSitemap\ SƠ ĐỒ TƯ DUY: MÔ HÌNH TIA SỐ TỰ NHIÊN \& BẢN CHẤT LIÊN TIẾP},
  colback=yellow!2!white,
  colframe=orange!85!black,
  arc=2.5mm, boxrule=0.8pt, top=1.5mm, bottom=1.5mm, left=1mm, right=1mm,
  halign=center
]
  \begin{tikzpicture}[
    xscale=1.05, yscale=0.82,
    tag/.style={rectangle, rounded corners=1.5mm, draw=teal!80!black, fill=teal!10, inner sep=2.5pt, font=\scriptsize\bfseries, align=center, text width=3.6cm},
    arr/.style={-{Stealth[scale=0.8]}, thick, draw=blue!80!black}
  ]
    % Tia số
    \draw[-{Stealth[scale=1.2]}, very thick, draw=teal!85!black] (-0.3, 0) -- (12.4, 0);
    
    \foreach \x/\val in {0/0, 1.4/1, 2.8/2, 4.2/3, 5.6/4, 9.2/n, 10.8/{n+1}} {
      \draw[thick, draw=teal!90!black] (\x, 0.15) -- (\x, -0.15);
      \node[below=2pt, font=\small\bfseries] at (\x, -0.15) {$\val$};
    }
    \node[below=3pt, font=\small\bfseries] at (7.4, -0.15) {$\dots$};

    \node[above=8pt, font=\scriptsize\bfseries\color{red!80!black}] at (0, 0.4) {Gốc 0 (Bé nhất)};
    \draw[->, thick, draw=red!75!black] (0, 0.4) -- (0, 0.18);
    \node[right=2pt, font=\scriptsize\bfseries\color{teal!85!black}] at (12.4, 0) {Vô tận};

    % Cung nhảy +1
    \draw[arr] (1.4, 0.35) to[bend left=35] node[above, font=\tiny\bfseries]{$+1$} (2.8, 0.35);
    \draw[arr] (2.8, 0.35) to[bend left=35] node[above, font=\tiny\bfseries]{$+1$} (4.2, 0.35);
    \draw[arr] (4.2, 0.35) to[bend left=35] node[above, font=\tiny\bfseries]{$+1$} (5.6, 0.35);
    \draw[arr] (9.2, 0.35) to[bend left=35] node[above, font=\tiny\bfseries]{$+1$} (10.8, 0.35);

    % Hộp ghi chú phía dưới
    \node[tag, fill=blue!10, draw=blue!70!black] at (2.0, -1.25) {Số chẵn tận cùng:\\$0, 2, 4, 6, 8$};
    \node[tag, fill=orange!12, draw=orange!75!black] at (6.2, -1.25) {Hai số liên tiếp:\\Hơn kém nhau $1$ đơn vị};
    \node[tag, fill=purple!10, draw=purple!70!black] at (10.4, -1.25) {Số lẻ tận cùng:\\$1, 3, 5, 7, 9$};
  \end{tikzpicture}
\end{tcolorbox}

\vspace{2mm}

% --- KHỐI 4: VÍ DỤ MINH HỌA & BẪY TRẮC NGHIỆM ---
\noindent
\begin{minipage}[t]{0.53\textwidth}
  \begin{tcolorbox}[
    enhanced,
    title={\footnotesize\bfseries\color{white}\faCheckCircle\ Ví dụ thực hành chuẩn mực},
    colback=green!2!white,
    colframe=green!60!black,
    arc=2.5mm, boxrule=0.8pt, top=1.8mm, bottom=1.8mm, left=2.2mm, right=2.2mm
  ]
    \footnotesize
    $\bullet$ \textbf{Ví dụ 1:} Viết 3 số tự nhiên liên tiếp biết số ở giữa là $999$:\\
    Số liền trước: $999 - 1 = 998$. Số liền sau: $999 + 1 = 1000$.\\
    Ba số liên tiếp là: $\mathbf{998,\, 999,\, 1000}$.\\
    $\bullet$ \textbf{Ví dụ 2:} Viết 3 số chẵn liên tiếp bắt đầu từ $2024$:\\
    Vì hai số chẵn liên tiếp cách nhau 2 đơn vị nên ta cộng tiếp 2:\\
    $\mathbf{2024,\, 2026,\, 2028}$.
  \end{tcolorbox}
\end{minipage}%
\hfill
\begin{minipage}[t]{0.45\textwidth}
  \begin{tcolorbox}[
    enhanced,
    title={\footnotesize\bfseries\color{white}\faExclamationTriangle\ Mẹo tránh bẫy trắc nghiệm},
    colback=red!2!white,
    colframe=red!70!black,
    arc=2.5mm, boxrule=0.8pt, top=1.8mm, bottom=1.8mm, left=2.2mm, right=2.2mm
  ]
    \footnotesize
    $\bullet$ \textbf{Bẫy số tự nhiên nhỏ nhất:}\\
    Nhiều bạn ghi số tự nhiên nhỏ nhất là $1$ $\rightarrow$ \textbf{SAI}! Số tự nhiên nhỏ nhất là số \textbf{0}.\\
    $\bullet$ \textbf{Bẫy số liền trước số 0:}\\
    Trong tập số tự nhiên, số 0 \textbf{không có số liền trước}.\\
    $\bullet$ \textbf{Bẫy liên tiếp}: Phải đọc kĩ là "liên tiếp" ($+1$) hay "chẵn/lẻ liên tiếp" ($+2$).
  \end{tcolorbox}
\end{minipage}

\end{document}
